\documentclass[aspectratio=169]{beamer}
\usepackage{tikz}
\usepackage{graphicx}
\usepackage{listings}
\usepackage{xcolor}
\usepackage{hyperref}
\usepackage[spanish,es-noquoting]{babel}
\usepackage[round]{natbib}

\lstset{
  basicstyle=\ttfamily\small,
  keywordstyle=\color{blue},
  stringstyle=\color{red},
  commentstyle=\color{gray},
  showstringspaces=false
}

% Clean Theme Setup & Optimization
\usetheme{Warsaw}
\usecolortheme{default}
\usefonttheme{serif}

% Remove top headline bar and navigation clutter to maximize slide text space
\setbeamertemplate{headline}{}
\setbeamertemplate{navigation symbols}{}

% Slim, clean footline bar
\setbeamertemplate{footline}{
  \leavevmode%
  \hbox{%
    \begin{beamercolorbox}[wd=.333333\paperwidth,ht=2.25ex,dp=1ex,left]{author in head/foot}%
      \hspace*{2ex}\insertshortauthor
    \end{beamercolorbox}%
    \begin{beamercolorbox}[wd=.333333\paperwidth,ht=2.25ex,dp=1ex,center]{title in head/foot}%
      \insertshorttitle
    \end{beamercolorbox}%
    \begin{beamercolorbox}[wd=.333333\paperwidth,ht=2.25ex,dp=1ex,right]{date in head/foot}%
      \insertframenumber{} / \inserttotalframenumber\hspace*{2ex}
    \end{beamercolorbox}%
  }%
  \vskip0pt%
}

% Font size and list spacing optimizations to prevent text overflow
\setbeamertemplate{itemize/enumerate body begin}{\small}
\setbeamertemplate{itemize/enumerate subbody begin}{\footnotesize}
\setbeamertemplate{itemize/enumerate subsubbody begin}{\scriptsize}

\makeatletter
\g@addto@macro\itemize{%
  \setlength{\itemsep}{2pt}%
  \setlength{\parsep}{1pt}%
  \setlength{\topsep}{2pt}%
  \setlength{\partopsep}{0pt}%
}
\makeatother

\graphicspath{{../}{./}}
\setbeamertemplate{caption}[numbered]

% Title info
\title[Kyverno: Vigilancia y Ejecució]{Kyverno: Vigilancia y Ejecución de Políticas de Seguridad en Kubernetes}
\subtitle{Una Perspectiva del Oficial de Políticas de Seguridad}
\author{Oficial Superior de Políticas de Seguridad}
\institute{Academia de Ciberseguridad / Cuerpo de Seguridad del Estado}
\date{\today}



\begin{document}

% Title slide
\begin{frame}
  \titlepage
\end{frame}

% Reset background for normal slides
\usebackgroundtemplate{}


\begin{frame}[allowframebreaks,fragile]{Kyverno: Vigilancia y Ejecución de Políticas de Seguridad en Kubernetes}

\textbf{Una Perspectiva del Oficial de Políticas de Seguridad}

\end{frame}

\begin{frame}[allowframebreaks,fragile]{Objetivos de la Sesión y Contexto General}


\textbf{Objetivos}
\begin{itemize}
  \item \textbf{Comprender} el rol fundamental de Kyverno en la seguridad de entornos Kubernetes.
  \item \textbf{Identificar} las capacidades clave de Kyverno para la aplicación efectiva de políticas.
  \item \textbf{Analizar} su impacto en el fortalecimiento de la postura de seguridad y el cumplimiento normativo.
  \item \textbf{Explorar} casos de uso prácticos y su integración en estrategias de DevSecOps y respuesta a incidentes.
\end{itemize}
\textbf{Contexto General}
En el panorama tecnológico actual, las arquitecturas nativas de la nube y los clústeres de Kubernetes son omnipresentes. Sin embargo, su complejidad inherente introduce desafíos de seguridad significativos. La necesidad de automatizar la aplicación de políticas de seguridad, garantizar el cumplimiento normativo y mitigar riesgos de forma proactiva es fundamental para cualquier organización moderna. Herramientas como Kyverno se vuelven esenciales para mantener una postura de seguridad robusta y adaptativa.

\end{frame}

\begin{frame}[allowframebreaks,fragile]{Marco Conceptual y Conceptos Clave: ¿Qué es Kyverno?}


Kyverno es un motor de políticas de código abierto diseñado específicamente para Kubernetes. Permite gestionar, validar, mutar y generar recursos de Kubernetes basándose en políticas definidas, actuando directamente dentro del clúster. Esto lo convierte en un componente crucial para la gobernanza y la seguridad en entornos cloud-native.

\textbf{Capacidades Fundamentales de Kyverno}
\begin{itemize}
  \item \textbf{Validación:} Asegura que los recursos de Kubernetes cumplen con reglas de seguridad predefinidas \textit{antes} de ser admitidos en el clúster.
  \begin{itemize}
    \item \textit{Ejemplo práctico:} Prevenir el despliegue de imágenes de contenedor que no provengan de un registro de imágenes aprobado o que carezcan de una etiqueta de seguridad obligatoria.
  \end{itemize}
  \item \textbf{Mutación:} Modifica automáticamente los recursos para aplicar configuraciones de seguridad estándar o corregir incumplimientos de políticas de forma proactiva.
  \begin{itemize}
    \item \textit{Ejemplo práctico:} Añadir automáticamente un \texttt{securityContext} a todos los Pods para forzar la ejecución con un usuario no-root o establecer límites de recursos (\texttt{CPU}, \texttt{memoria}) para prevenir abusos.
  \end{itemize}
\end{itemize}
\end{frame}

\begin{frame}[allowframebreaks,fragile]{Marco Conceptual y Conceptos Clave: Funcionalidades Avanzadas y Relevancia}


\textbf{Capacidades Fundamentales de Kyverno (Continuación)}
\begin{itemize}
  \item \textbf{Generación:} Crea nuevos recursos de Kubernetes basándose en políticas, facilitando la implementación de configuraciones seguras por defecto o la estandarización de recursos de seguridad.
  \begin{itemize}
    \item \textit{Ejemplo práctico:} Generar automáticamente políticas de red (\texttt{NetworkPolicies}) o roles RBAC para nuevos espacios de nombres (\texttt{namespaces}) creados, asegurando un aislamiento y control de acceso consistentes desde el inicio.
  \end{itemize}
\end{itemize}
\textbf{La Perspectiva del Oficial de Políticas de Seguridad}
Kyverno se posiciona como una herramienta indispensable para:
\begin{itemize}
  \item \textbf{Aplicación Consistente de Políticas de Seguridad:} Garantiza que las configuraciones y despliegues en Kubernetes se adhieran rigurosamente a los requisitos de seguridad definidos por la organización.
  \item \textbf{Cumplimiento Normativo Automatizado:} Facilita la observancia de estándares y regulaciones sectoriales (ej., ISO 27001, PCI DSS, GDPR) al automatizar la verificación y el cumplimiento de las políticas.
  \item \textbf{Mitigación Proactiva de Riesgos:} Reduce drásticamente la superficie de ataque al prevenir despliegues inseguros, corregir configuraciones erróneas y responder a amenazas en tiempo real.
\end{itemize}
\end{frame}

\begin{frame}[allowframebreaks,fragile]{Marco Conceptual y Conceptos Clave: Política como Código y DevSecOps}


\textbf{Kyverno y el Enfoque de ``Política como Código'' (Policy-as-Code)}
\begin{itemize}
  \item \textbf{Definición Transparente:} Las políticas en Kyverno se definen como recursos de Kubernetes (Custom Resources, CRDs), lo que permite tratarlas como código.
  \item \textbf{Beneficios Clave:}
  \begin{itemize}
    \item \textbf{Control de Versiones (GitOps):} Gestionar políticas en sistemas como Git, permitiendo la trazabilidad, reversión y auditoría de cambios.
    \item \textbf{Pruebas Automatizadas:} Integrar pruebas de políticas en el CI/CD para validar su comportamiento antes del despliegue.
    \item \textbf{Consistencia y Escalabilidad:} Asegurar la aplicación uniforme de políticas en múltiples clústeres y entornos.
    \item \textbf{Mejora de la Auditabilidad:} Simplificar los procesos de auditoría al tener las políticas versionadas y documentadas junto con el resto de la infraestructura.
  \end{itemize}
\end{itemize}
\textbf{Integración en Estrategias de DevSecOps}
\begin{itemize}
  \item Kyverno permite integrar controles de seguridad temprana en el ciclo de vida del desarrollo y despliegue de software, automatizando las comprobaciones en tiempo de admisión. Esto ayuda a `\texttt{desplazar la seguridad a la izquierda'' (}shift left`), identificando y corrigiendo vulnerabilidades antes de que lleguen a producción.
\end{itemize}
\end{frame}

\begin{frame}[allowframebreaks,fragile]{Casos de Estudio y Ejemplos Prácticos Aplicados}


\textbf{1. Integridad de la Cadena de Suministro de Software}
\begin{itemize}
  \item \textbf{Referencia:} Guduru, S. (2021). \textit{Automated Policy Enforcement in DevSecOps: OPA, Kyverno, and Tekton Chains for Supply Chain Integrity.} International Journal of Science and Research (IJSR). DOI: 10.21275/sr210513123637
  \item \textbf{Aplicación Práctica para el Oficial de Seguridad:} Kyverno puede hacer cumplir políticas que validen la integridad y el origen de los artefactos de software.
  \begin{itemize}
    \item \textbf{Ejemplo:} Una política de Kyverno podría requerir que todas las imágenes de contenedor desplegadas en el clúster estén firmadas digitalmente por una clave de confianza específica y provengan exclusivamente de un registro de imágenes interno aprobado. Esto previene el uso de imágenes comprometidas o no autorizadas, fortaleciendo la seguridad de la cadena de suministro.
  \end{itemize}
\end{itemize}
\end{frame}

\begin{frame}[allowframebreaks,fragile]{Casos de Estudio y Ejemplos Prácticos Aplicados}


\textbf{2. Mitigación de Amenazas en Tiempo de Ejecución con Respuesta Automatizada}
\begin{itemize}
  \item \textbf{Referencia:} Shenoy, A., S. P, Honnavali, P. (2025). \textit{Runtime Threat Mitigation in Kubernetes Using Falco, Falcosidekick, and Kyverno: An Automated Pod Deletion Approach.} 2025 IEEE 7th International Conference on Computing, Communication and Automation (ICCCA). DOI: 10.1109/iccca66364.2025.11325363
  \item \textbf{Aplicación Práctica para el Oficial de Seguridad:} Kyverno, en conjunto con herramientas de detección de amenazas en tiempo real como Falco, permite una respuesta automatizada a incidentes de seguridad.
  \begin{itemize}
    \item \textbf{Ejemplo:} Si Falco detecta un comportamiento anómalo y malicioso dentro de un Pod (ej., ejecución de un shell inesperado, intentos de escalada de privilegios), una política de Kyverno podría configurarse para eliminar automáticamente ese Pod comprometido. Esta acción rápida reduce la ventana de oportunidad para un atacante y minimiza el impacto de la brecha.
  \end{itemize}
\end{itemize}
\textbf{3. Defensa Multi-Capa contra Amenazas Específicas: Minería de Criptomonedas}
\begin{itemize}
  \item \textbf{Referencia:} Sagaram, T., S. P, Honnavalli, P. (2025). \textit{Multi-Layered Runtime Defense Against Cryptocurrency Mining in Kubernetes Using Falco, Kyverno, and Behavior Analytics.} 2025 66th International Scientific Conference on Information Technology and Management Science of Riga Technical University (ITMS). DOI: 10.1109/itms67030.2025.11236625
  \item \textbf{Aplicación Práctica para el Oficial de Seguridad:} Kyverno puede reforzar una estrategia de defensa multi-capa contra amenazas específicas como el criptominado no autorizado.
  \begin{itemize}
    \item \textbf{Ejemplo:} Kyverno puede aplicar políticas de mutación para establecer límites estrictos de CPU y memoria en todos los Pods, o generar políticas de red que prohíban la comunicación con puertos y destinos asociados a la minería de criptomonedas. Esto protege los recursos del clúster y previene pérdidas económicas y operativas.
  \end{itemize}
\end{itemize}
\end{frame}

\begin{frame}[allowframebreaks,fragile]{Ejercicios Prácticos y Actividades Interactivas}


\textbf{Ejercicio 1: Identificación de Violaciones y Diseño de Políticas de Validación (5 minutos)}
\textbf{Escenario:} Su organización tiene una política que prohíbe el despliegue de contenedores con privilegios de \texttt{root} y requiere que todos los Pods tengan límites de recursos definidos (\texttt{CPU} y \texttt{memoria}) para evitar un consumo excesivo.

\begin{itemize}
  \item \textbf{Identifique la Violación:} Un desarrollador intenta desplegar un Pod sin límites de recursos y con \texttt{privileged: true} en su \texttt{securityContext}. Identifique las violaciones de políticas de seguridad.
  \item \textbf{Diseñe una Política de Validación con Kyverno (Conceptual):} Describa cómo una política de Kyverno podría \textbf{validar} este Pod, impidiendo su creación si no cumple con estas dos reglas.
  \begin{itemize}
    \item \textit{Pistas:} Piense en las propiedades \texttt{securityContext.privileged} y \texttt{resources.limits} dentro de la especificación del Pod.
  \end{itemize}
\end{itemize}
\end{frame}

\begin{frame}[allowframebreaks,fragile]{Ejercicios Prácticos y Actividades Interactivas}


\textbf{Ejercicio 2: Reforzando el Cumplimiento con Políticas de Mutación (5 minutos)}
\textbf{Escenario:} Una nueva política de seguridad exige que todos los \texttt{Deployments} recién creados en el clúster tengan configurados volúmenes de \texttt{emptyDir} con el modo de acceso \texttt{readOnlyRootFilesystem: true} para mejorar la seguridad de los archivos temporales.

\begin{itemize}
  \item \textbf{Problema Actual:} Un nuevo \texttt{Deployment} se crea con un volumen \texttt{emptyDir} que no tiene \texttt{readOnlyRootFilesystem: true}.
  \item \textbf{Proponga una Política de Mutación con Kyverno (Conceptual):} ¿Cómo configuraría una política de Kyverno para que automáticamente \textbf{mute} cualquier nuevo \texttt{Deployment} que incluya un \texttt{emptyDir} y le añada la propiedad \texttt{readOnlyRootFilesystem: true} si está ausente?
  \begin{itemize}
    \item \textit{Pistas:} Considere la ruta del recurso (e.g., \texttt{spec.template.spec.volumes}) y las propiedades específicas del volumen.
  \end{itemize}
\end{itemize}
\end{frame}

\begin{frame}[allowframebreaks,fragile]{``Sabías que...'' / Datos Curiosos}


\begin{itemize}
  \item \textbf{Integración Nativa que Simplifica lo Complejo:} A diferencia de muchos motores de políticas que actúan como componentes externos, Kyverno se integra de forma nativa con Kubernetes, utilizando sus propios recursos (\texttt{Custom Resources}). Esto es sorprendente porque simplifica enormemente la definición y aplicación de políticas de seguridad complejas, haciendo que la gestión de la seguridad se sienta como una parte inherente del propio Kubernetes, más accesible y menos propensa a errores para operadores y equipos de seguridad.
  \item \textbf{Más allá de la Prevención: Remedio y Provisión Segura Automatizada:} A menudo se asocia la aplicación de políticas con la validación (es decir, el bloqueo de recursos no conformes). Sin embargo, Kyverno va un paso más allá con sus capacidades de mutación y generación. Esto significa que no solo puede bloquear un recurso inseguro, sino que puede modificarlo automáticamente para hacerlo seguro (ej., añadir un contexto de seguridad) o incluso crear nuevos recursos seguros (ej., políticas de red) de forma proactiva. Esta capacidad transforma la seguridad de una función puramente reactiva a una proactiva y automatizada.
  \item \textbf{Respuesta Dinámica en Tiempo de Ejecución y Amenazas Específicas:} La investigación académica destaca la capacidad de Kyverno para la seguridad dinámica en tiempo de ejecución, especialmente cuando se integra con herramientas de detección de amenazas como Falco. El hecho de que Kyverno pueda aplicar políticas basadas en el comportamiento observado en tiempo real, desencadenando respuestas automatizadas como la eliminación de Pods comprometidos o el aislamiento de recursos, es una revelación que expande su utilidad más allá de las simples comprobaciones de configuración estáticas. Esto posiciona a Kyverno como un componente crítico en una estrategia de defensa en profundidad contra vectores de ataque evolucionados, como la minería de criptomonedas no autorizada.
\end{itemize}
\end{frame}

\begin{frame}[allowframebreaks,fragile]{Conclusiones y Puntos Clave}


\begin{itemize}
  \item \textbf{Kyverno es un Guardián Esencial en Kubernetes:} Actúa como un motor de políticas robusto, garantizando que los recursos cumplan con las directrices de seguridad desde el momento de su admisión.
  \item \textbf{Automatización Integral de la Seguridad:} Facilita la aplicación de políticas de seguridad, cumplimiento y las mejores prácticas mediante la validación, mutación y generación de recursos de forma automatizada.
  \item \textbf{Impulsor del DevSecOps y la Integridad de la Cadena de Suministro:} Permite la integración temprana de la seguridad en el ciclo de vida del desarrollo y despliegue, reforzando la integridad de los artefactos de software.
  \item \textbf{Defensa Dinámica y Proactiva:} Va más allá de la prevención estática, permitiendo respuestas automatizadas y dinámicas a amenazas en tiempo de ejecución en conjunto con herramientas de detección.
  \item \textbf{Auditabilidad Mejorada y Consistencia Inquebrantable:} Su enfoque de ``política como código'' (Policy-as-Code) mejora la auditabilidad, la trazabilidad y la aplicación consistente de políticas en todos los entornos de Kubernetes.
\end{itemize}
Como Oficial de Políticas de Seguridad, la implementación estratégica de Kyverno es fundamental para mantener una postura de seguridad robusta, ágil y conforme en los complejos y dinámicos entornos cloud-native.

\end{frame}

\begin{frame}[allowframebreaks,fragile]{Referencias Bibliográficas}


\begin{itemize}
  \item Guduru, S. (2021). Automated Policy Enforcement in DevSecOps: OPA, Kyverno, and Tekton Chains for Supply Chain Integrity. \textit{International Journal of Science and Research (IJSR)}, 10(5). DOI: 10.21275/sr210513123637
  \item Shenoy, A., S. P, Honnavali, P. (2025). Runtime Threat Mitigation in Kubernetes Using Falco, Falcosidekick, and Kyverno: An Automated Pod Deletion Approach. \textit{2025 IEEE 7th International Conference on Computing, Communication and Automation (ICCCA)}. DOI: 10.1109/iccca66364.2025.11325363
  \item Sagaram, T., S. P, Honnavalli, P. (2025). Multi-Layered Runtime Defense Against Cryptocurrency Mining in Kubernetes Using Falco, Kyverno, and Behavior Analytics. \textit{2025 66th International Scientific Conference on Information Technology and Management Science of Riga Technical University (ITMS)}. DOI: 10.1109/itms67030.2025.11236625
\end{itemize}
\end{frame}




\end{document}